\section{Analysis}
\begin{figure}[H]
    \centering
    \includegraphics[width = \textwidth]{runnig_mean_4x4.png}
    \caption{Running mean of Plaquette for 4x4-Lattice.}
    \label{fig:run_mean_4_L}
\end{figure}

The equilibration of the simulated system is a crucial first step in numerical lattice simulations. 
\Cref{fig:run_mean_4_L} shows the running mean of the plaquette.
The plaquette $P_{\mu\nu}(n)$ measures the product of link variables around an elementary square on the lattice:
\begin{equation}
    P_{\mu\nu}(n) = \text{Re}[U_\mu(n)U_\nu(n+\mu)U_\mu^\dagger(n+\nu)U_\nu^\dagger(n)]
\end{equation}

The running mean demonstrates that the system rapidly converges to its equilibrium value within the first few hundred trajectories. 
The initial transient behavior, where the plaquette value drifts from its starting configuration, indicates that the system is thermalizing. 
Once the running mean stabilizes, the system has reached thermal equilibrium and subsequent measurements represent independent samples from the desired distribution. 
Similar observation can be seen for different lattice size \cref{fig:run_mean_5_L} and \cref{fig:run_mean_6_L}

\subsection{Topological Charge Evolution}
The integer-valued topological charge functional for 2D- Schwinger model is given by~\cite{elser1997topological},
\begin{equation}
Q\;=\; \frac{1}{2\pi}\int d^{2}x\;F_{01}(x)\,.
\end{equation}
One can measure the discreet topological charge with
\begin{equation} \label{eqn:discreet_Q}
Q=\frac{1}{2\pi}\sum_{n \in \Lambda}\arg P_{\mu\nu}(n)
\end{equation}
\begin{figure}[H]
    \centering
    \includegraphics[width=\textwidth]{Q_4.png}
    \caption{Topological charge for 4x4-Lattice}
    \label{fig:Q_4_L}
\end{figure}

\Cref{fig:Q_4_L}, \cref{fig:Q_5_L} and \cref{fig:Q_6_L} displays the distribution of the topological charge from the HMC simulation. 
One can observe that, the topological charge fluctuates between different integer values, 
demonstrating that the HMC algorithm successfully tunnels between distinct topological sectors.

This behavior is highly desirable and contrasts sharply with some local update algorithms that become 
frozen or have very low tunnelling probability in a single topological sector~\cite{elser1997topological}. 
The ability to explore multiple topological sectors ensures that:
\begin{itemize}
    \item The simulation samples from the complete configuration space
    \item Physical observables are not biased by topological sector restrictions
\end{itemize}

\subsection{Autocorrelation Analysis}

\begin{figure}
    \hspace*{-4cm}
    \includegraphics[width=1.6\textwidth]{autocorr_5_Lattice.png}
    \caption{Autocorrelation of topological charge (left) and Wilson loop W(R,T) for R = 1, T = 1 in 5x5-Lattice (right)}
\end{figure}

One of the primary advantages of the Hybrid Monte Carlo algorithm over local update methods is the significantly reduced autocorrelation time. 
\Cref{fig:Q_4_L} shows normalized autocorrelation function $\Gamma$ for topological charge  on left  and Wilson loop $W(R,T)$  on right. Both shows 
significant low integrated autocorrelation time. For topological charge $\tau_{\text{int}} \sim 1.28$ and for Wilson loop $W(1,1)$ $\tau_{\text{int}} \sim 0.5$ 
for 5x5-Lattice. Both the topological charge and the Wilson loop demonstrate
the similar autocorrelation and integrated autocorrelation time for different lattice size.

\subsection{Confinement Evidence}

\begin{figure}
    \centering
    \includegraphics[width=\textwidth]{pot_V_T_1.png}
    \caption{Potential plotted as -ln(W(R,T)) with T = 1. The potential growes linearly showing confinement.}
    \label{fig:pot_V_T_1}
\end{figure}

A defining characteristic of QCD and 2D-Schwinger model is the phenomenon of confinement. 
The confinement in Schwinger model is due to Gauss' law in 1-Spatial dimension $\partial_x E \sim q$.
Unlike free theories where the potential between charges grows logarithmically (or vanishes), confining theories exhibit a linear potential at large distances. 
In \cref{fig:pot_V_T_1} a linear fit to \cref{eq:V_R} demonstrates this confinement property by plotting the effective potential,

\begin{equation}
    \label{eq:V_R}
    V(R) = -\ln W(R,T) / T
\end{equation}

where $W(R,T)$ is the Wilson loop operator for a rectangular closed loop with spatial extent $R$ and temporal extent $T$~\cite{lecturenote}. 
The linear growth of the potential with separation distance $R$ shows confinement predicted by the theory.

This linear confinement behavior emerges naturally from the HMC simulations without any external constraint and is a consequence of the gauge-invariant 
dynamics of the Schwinger model. The successful observation of confinement demonstrates that:
\begin{itemize}
    \item The lattice formulation correctly preserves the essential physics
    \item The HMC algorithm efficiently samples the relevant gauge configurations
    \item The numerical implementation accurately reproduces theoretical predictions
\end{itemize}

The confinement phenomenon observed in these simulations validates the approach and confirms that the 2D Schwinger model, 
studied via lattice Monte Carlo methods, provides valuable insights for understanding gauge theories.

\newpage

\section{Conclusions}

The numerical simulations of the 2D Schwinger model using the Hybrid Monte Carlo algorithm have successfully demonstrated several key physical and computational results:

\begin{enumerate}
    \item \textbf{Rapid Equilibration:} The system reaches thermal equilibrium within a few hundred HMC trajectories, as evidenced by the stabilization of the plaquette running mean. This rapid convergence validates the choice of HMC for exploring the configuration space efficiently.
    
    \item \textbf{Topological Sector Exploration:} The algorithm successfully tunnels between different topological sectors, with the topological charge fluctuating between integer values throughout the simulation. This demonstrates that HMC provides unbiased sampling across the entire configuration manifold, avoiding the freezing that plagues local update algorithms.
    
    \item \textbf{Reduced Autocorrelation:} Both the topological charge and Wilson loop observables exhibit short autocorrelation times, indicating statistically independent measurements. This efficiency gain over local methods translates to faster convergence of statistical estimates and reduced computational cost for achieving target precision.
    
    \item \textbf{Confinement Behavior:} The extracted interquark potential demonstrates clear linear growth with distance, confirming the confinement property expected from Schwinger model. 
\end{enumerate}


Despite the successful demonstration of HMC's capabilities, the simulations also reveal important limitations that become increasingly significant as the lattice volume increases.
The leapfrog integrator, introduces systematic errors that scale as $\mathcal{O}(\epsilon^2)$ per step. 
For a fixed trajectory length $\tau$, the leapfrog (second-order) integrator
induces a Hamiltonian violation whose typical size scales as
\begin{equation}
\Delta H_{\mathrm{rms}} \sim \sqrt{V}\,\epsilon^{2},
\qquad
\epsilon = \frac{\tau}{N_{\text{step}}}.
\end{equation}
Consequently, increasing the lattice volume $V$ leads to larger integration
errors and a reduced Metropolis acceptance probability.
To maintain an acceptance rate above $50\%$, the step size $\epsilon$ must be
decreased, which in practice requires increasing the number of leapfrog steps
$N_{\text{step}}$ at fixed $\tau$.
This significantly increases the computational cost, since each additional
molecular-dynamics step requires an evaluation of the  force.